%%%%%%%%%%%%%%%%%%%%%%%%%%%%%%%%%%%%%%%%%%%%%%%%%%%%%%%%%%%%%%%%%%%%%%%%%%%
%
% Generic template for TFC/TFM/TFG/Tesis
%
% By:
%  + Javier Macías-Guarasa.
%    Departamento de Electrónica
%    Universidad de Alcalá
%  + Roberto Barra-Chicote.
%    Departamento de Ingeniería Electrónica
%    Universidad Politécnica de Madrid
%
% By: + Javier Macías-Guarasa. Departamento de Electrónica Universidad de Alcalá + Roberto Barra-Chicote. Departamento de Ingeniería Electrónica Universidad Politécnica de Madrid
%
% Based on original sources by Roberto Barra, Manuel Ocaña, Jesús Nuevo, Pedro Revenga, Fernando Herránz and Noelia Hernández. Thanks a lot to all of them, and to the many anonymous contributors found (thanks to google) that provided help in setting all this up.
%
% See also the additionalContributors.txt file to check the name of additional contributors to this work.
%
% If you think you can add pieces of relevant/useful examples, improvements, please contact us at (macias@depeca.uah.es)
%
% You can freely use this template and please contribute with comments or suggestions!!!
%
%%%%%%%%%%%%%%%%%%%%%%%%%%%%%%%%%%%%%%%%%%%%%%%%%%%%%%%%%%%%%%%%%%%%%%%%%%%

\chapter{Compilación y herramientas avanzadas}
\label{cha:compilacion-avanzada}

No necesitas las herramientas de este capítulo para comenzar a escribir o para compilar el documento desde tu editor. Las hemos reunido aquí para que puedas automatizar tareas, preparar revisiones y diagnosticar problemas cuando tu forma de trabajo las requiera.

\section{Compilación mediante el \texttt{Makefile}}
\label{sec:compilacion-makefile}

Para facilitar la generación del documento incluimos un \texttt{Makefile}. No somos expertos ni en \texttt{make} ni en \LaTeX{}, así que seguramente no sea el mejor \texttt{Makefile} del mundo, pero creemos que cumple su función y puede servir también como ejemplo para adaptar el proceso de compilación a tus necesidades.

El objetivo predeterminado \texttt{all}, ejecutado al invocar \texttt{make} sin argumentos, convierte primero los diagramas y figuras que lo necesiten. A continuación ejecuta \texttt{pdflatex}, Biber, \texttt{makeglossaries} y las dos pasadas finales de \texttt{pdflatex}. Dado que el libro de ejemplo utiliza acrónimos y símbolos, el \texttt{Makefile} ejecuta siempre \texttt{makeglossaries}.

La compilación genera \texttt{book.pdf} y \texttt{book-compressed.pdf}. También crea copias cuyos nombres identifican el tipo de trabajo, la titulación, el autor y el idioma configurados. La versión comprimida utiliza la configuración \texttt{/prepress} de Ghostscript; la antigua versión de calidad \texttt{/screen} está desactivada porque degradaba excesivamente las imágenes.

El objetivo \texttt{all} también actualiza el fichero \texttt{RELEASE.txt} con la etiqueta más reciente del repositorio e intenta añadirlo al índice de Git. Por tanto, después de compilar, revisa el estado del repositorio antes de preparar un commit.

Para el uso habitual de la plantilla solo necesitarás los siguientes objetivos:

\begin{itemize}
  \item \texttt{all}: compila el documento completo y genera las copias finales; es el objetivo ejecutado de forma predeterminada por \texttt{make}.
  \item \texttt{clean}: elimina los ficheros auxiliares y los PDF de trabajo generados durante la compilación.
\end{itemize}

Los siguientes objetivos proporcionan operaciones avanzadas para la revisión, la preparación o la distribución del proyecto:

\begin{itemize}
  \item \texttt{all\_latexmk}: ofrece un procedimiento alternativo y experimental de compilación basado en \texttt{latexmk}.
  \item \texttt{clean\_latexmk}: elimina los ficheros auxiliares mediante \texttt{latexmk} y completa después la limpieza de otros ficheros generados.
  \item \texttt{snapshot}: genera mediante \texttt{latexpand} una copia unificada del documento que puede utilizarse como referencia para una revisión posterior.
  \item \texttt{flatten}: genera una copia unificada del estado actual del documento.
  \item \texttt{latexdiff}: compara la instantánea con el estado actual y genera un PDF que muestra las diferencias.
  \item \texttt{backup-chapters}: crea copias fechadas de los capítulos y apéndices.
  \item \texttt{orig-chapters}: sustituye los capítulos y apéndices de trabajo por las versiones completas de ejemplo, después de crear una copia de seguridad.
  \item \texttt{bare-chapters}: sustituye los capítulos y apéndices de trabajo por versiones mínimas, después de crear una copia de seguridad.
  \item \texttt{bare}: prepara una estructura mínima mediante \texttt{bare-chapters} y ofrece eliminar los gráficos y diagramas existentes.
  \item \texttt{tar}: genera archivos comprimidos con los fuentes y recursos de la plantilla.
\end{itemize}

Los objetivos relacionados con la comparación de versiones se explican con más detalle en la sección~\ref{sec:control-de-cambios}.

El \texttt{Makefile} contiene además objetivos auxiliares utilizados internamente para convertir recursos y preparar la versión mínima. No necesitas ejecutarlos directamente.

Los objetivos \texttt{bare}, \texttt{bare-chapters} y \texttt{orig-chapters} modifican el contenido de los directorios de capítulos o apéndices. Además, \texttt{bare} puede eliminar los ficheros de \texttt{figures/} y \texttt{diagrams/} si se confirma expresamente durante su ejecución. Utilízalos únicamente al preparar una copia nueva del proyecto y después de comprobar que los cambios importantes están guardados en el sistema de control de versiones.

Algunos objetivos del \texttt{Makefile} utilizan herramientas adicionales. Solo necesitas instalar las que correspondan a los objetivos que vayas a utilizar:

\begin{itemize}
  \item \texttt{Ghostscript}: genera la copia comprimida del documento cuando se utiliza el objetivo \texttt{all} del \texttt{Makefile}.
  \item \texttt{Inkscape}: convierte automáticamente los diagramas en formato \texttt{svg} a \texttt{pdf}.
  \item \texttt{Dia} y \texttt{epspdf}: convierten automáticamente los diagramas en formato \texttt{dia} a \texttt{pdf}; \texttt{epspdf} también convierte las figuras en formato \texttt{eps} a \texttt{pdf}.
  \item \texttt{latexmk}: proporciona el procedimiento alternativo utilizado por los objetivos \texttt{all\_latexmk} y \texttt{clean\_latexmk}.
  \item \texttt{latexpand} y \texttt{latexdiff}: permiten unificar los fuentes y generar documentos que muestran las diferencias entre dos versiones.
  \item \texttt{tar}, \texttt{gzip} y \texttt{zip}: generan los archivos comprimidos creados por el objetivo \texttt{tar}.
  \item \texttt{Perl}: es necesario para ejecutar \texttt{makeglossaries} en los sistemas que no incluyan ya un intérprete compatible.
\end{itemize}

Instala estas herramientas solo cuando vayas a utilizar la prestación correspondiente. No son requisitos para redactar ni compilar un documento básico que no dependa de ellas.


\section{Generación del documento con control de cambios}
\label{sec:control-de-cambios}

Uno de los problemas que encontramos al trabajar con \LaTeX{} aparece durante la revisión: después de entregar un borrador, el tutor o director puede querer comprobar cómo se han incorporado sus sugerencias. En un procesador de textos como LibreOffice (vaaaaaaaaale, o Microsoft Word también) resulta habitual activar el control de cambios o comparar dos documentos; con \LaTeX{} también podemos obtener un resultado equivalente mediante algunas herramientas adicionales.

Para ello utilizamos \texttt{latexdiff}, que compara dos versiones de un documento \LaTeX{} y genera un nuevo fichero en el que se marcan visualmente las inserciones y eliminaciones. Al compilarlo obtenemos un PDF independiente mucho más cómodo para revisar que un \texttt{diff} a palo seco de los ficheros fuente.

Como el libro se divide mediante múltiples órdenes \texttt{\textbackslash{}input}, antes de realizar la comparación se utiliza \texttt{latexpand} para generar una versión unificada. El \texttt{Makefile} de \texttt{Book/} automatiza este procedimiento mediante los objetivos \texttt{snapshot}, \texttt{flatten} y \texttt{latexdiff}.

\subsection{Creación de la versión de referencia}

Cuando tengas una versión que quieras utilizar como base de la revisión, ejecuta desde el directorio \texttt{Book/}:

\begin{verbatim}
make snapshot
\end{verbatim}

Este objetivo genera \texttt{book-flatten-snapshot.tex}, que contiene en un solo fichero el estado expandido del documento. Crea la instantánea antes de comenzar las modificaciones que quieras mostrar y consérvala sin cambios durante el proceso de revisión.

Si utilizas Git, te recomendamos registrar la instantánea junto con la versión del documento a la que corresponde:

\begin{verbatim}
git add book-flatten-snapshot.tex
git commit -m "chore: update the review snapshot"
\end{verbatim}

El objetivo \texttt{snapshot} sobrescribe la instantánea existente. Por tanto, no vuelvas a ejecutarlo hasta que quieras establecer una nueva versión de referencia.

\subsection{Generación del documento de diferencias}

Después de modificar el documento, puedes generar la comparación mediante:

\begin{verbatim}
make latexdiff
\end{verbatim}

El objetivo \texttt{latexdiff} genera primero \texttt{book-flatten.tex} con el estado actual, lo compara con \texttt{book-flatten-snapshot.tex} y crea \texttt{book-flatten-diff.tex}. A continuación procesa su bibliografía y sus glosarios y genera \texttt{book-flatten-diff.pdf}, que contiene las diferencias resaltadas.

La comparación depende de que exista previamente \texttt{book-flatten-snapshot.tex}. Si quieres iniciar un nuevo ciclo de revisión, comprueba primero que hayas guardado los cambios anteriores y ejecuta de nuevo \texttt{make snapshot} para reemplazar la versión de referencia.

El \texttt{Makefile} aplica además una corrección específica al resultado de \texttt{latexdiff} para tratar determinados entornos de TikZ utilizados por la plantilla. Esta corrección ayuda con los ejemplos actuales, pero un documento que contenga construcciones complejas puede requerir ajustes manuales en \texttt{book-flatten-diff.tex}.

\subsection{Uso sin \texttt{make}}

El flujo de revisión no requiere obligatoriamente la utilidad \texttt{make}. Cristina Losada aportó un procedimiento para utilizar \texttt{latexdiff} directamente en Windows, que también puede adaptarse a otros sistemas operativos:

\begin{enumerate}
  \item Instala \texttt{latexdiff} y un intérprete de Perl. MiKTeX incluye \texttt{latexdiff}, y Strawberry Perl es una de las alternativas disponibles para Windows.
  \item Guarda una copia del documento que utilizarás como referencia; por ejemplo, \texttt{book-snapshot.tex}.
  \item Realiza las modificaciones necesarias en el documento.
  \item Ejecuta desde el directorio \texttt{Book/}:
\end{enumerate}

\begin{verbatim}
latexdiff --flatten book-snapshot.tex book.tex > book-flatten-diff.tex
\end{verbatim}

Después, compila \texttt{book-flatten-diff.tex}. Si la versión comparada contiene bibliografía, acrónimos o símbolos, ejecuta también Biber y \texttt{makeglossaries} con el nombre base \texttt{book-flatten-diff}, seguidos de las pasadas adicionales de \texttt{pdflatex}.

Este procedimiento manual compara directamente el documento guardado con la versión actual. Para que \texttt{--flatten} pueda expandir correctamente ambos documentos, conserva también los ficheros incluidos por cada versión o utiliza copias ya unificadas mediante \texttt{latexpand}.

También existen otras alternativas integradas con Git, como \texttt{git-latexdiff}~\cite{git-latexdiff}. El procedimiento incluido en este repositorio no depende de una herramienta concreta de control de versiones, aunque te recomendamos utilizar Git u otro sistema equivalente para conservar el historial de tu trabajo.

\section{Resolución de problemas}
\label{sec:problemas-conocidos}

Los problemas observados antiguamente con Ubuntu 12.04, el paquete \texttt{xcolor}, Evince y Xpdf ya no describen entornos actuales y se han retirado de esta guía. Si una compilación falla, localiza el primer error del registro, ya que muchos de los mensajes posteriores suelen ser consecuencias de ese fallo inicial.

Las comprobaciones más habituales son:

\begin{itemize}
  \item Si falta un fichero \texttt{.sty}, instala el paquete correspondiente desde la distribución de \LaTeX{} que utilices.
  \item Si no aparece la bibliografía o las citas quedan sin resolver, comprueba que hayas ejecutado Biber y realiza después las pasadas adicionales de \texttt{pdflatex}.
  \item Si no aparecen los listados de acrónimos o símbolos, ejecuta \texttt{makeglossaries} antes de las pasadas finales de \texttt{pdflatex}.
  \item Si falla la conversión de un fichero \texttt{dia}, \texttt{svg} o \texttt{eps}, instala la herramienta opcional correspondiente o proporciona directamente una versión compatible en \texttt{pdf}, \texttt{png} o \texttt{jpg}.
  \item Si las referencias cruzadas o los índices no están actualizados, repite las pasadas de \texttt{pdflatex}; antes de recompilar también puedes utilizar \texttt{make clean} para eliminar los auxiliares anteriores.
  \item Si \texttt{latexdiff} genera un fichero que no compila, revisa el primer error de \texttt{book-flatten-diff.tex}. Las tablas, fórmulas, entornos TikZ y órdenes personalizadas pueden necesitar correcciones manuales en el documento de diferencias.
  \item Si una cubierta o un formulario no coincide con el modelo institucional esperado, comprueba \texttt{\textbackslash{}myDegree}, el idioma, los datos compartidos y la normativa vigente antes de modificar directamente la plantilla.
\end{itemize}

Cuando nos comuniques un problema, incluye el sistema operativo, la distribución y versión de \LaTeX{}, la titulación seleccionada, el comando de compilación y el primer error relevante del fichero de registro. También puedes proponer correcciones mediante el repositorio del proyecto.

Si encuentras algún otro problema, escribe a \contactauthor{} para contárnoslo y trataremos de solucionarlo. Y si ya has encontrado una solución, no dudes en compartirla: tu contribución puede ser útil para quienes utilicen la plantilla después.


%%% Local Variables:
%%% TeX-master: "../book"
%%% End:
